
% !TeX program = lualatex
% !TeX encoding = utf8
% !TeX spellcheck = uk_UA
% !TeX root =../PyscfBook.tex

%=========================================================
\chapter{Розрахунки твердих тіл}\label{\currfilebase}
%=========================================================

%% --------------------------------------------------------
\section{Вступ: періодичні граничні умови}
%% --------------------------------------------------------

У молекулярних розрахунках ми маємо скінченну систему атомів, що ізольована у просторі.
Для твердих тіл ситуація інша --- структура нескінченно повторюється у просторі.
Щоб математично описати це, використовують \textbf{періодичні граничні умови} (Periodic Boundary Conditions, PBC).

Модуль \texttt{pyscf.pbc} дозволяє виконувати квантово-хімічні розрахунки для таких систем: кристалів, поверхонь, двовимірних матеріалів (як графен), наноструктур.
Тут основою є \textbf{теорія Блоха}, згідно з якою хвильові функції електронів у періодичному потенціалі записуються як:
\[
    \psi_{n\mathbf{k}}(\mathbf{r}) = e^{i \mathbf{k} \cdot \mathbf{r}} u_{n\mathbf{k}}(\mathbf{r}),
\]
де $\mathbf{k}$ --- вектор у зоні Бріллюена, $u_{n\mathbf{k}}(\mathbf{r})$ --- періодична частина, що повторюється з періодом ґратки.

\textbf{Основні відмінності від молекулярних розрахунків:}
\begin{itemize}
    \item Система описується не координатами всіх атомів, а \textbf{елементарною коміркою} та її векторами трансляції.
    \item Для зменшення кількості електронів використовують \textbf{псевдопотенціали}, що замінюють остовні електрони ефективним потенціалом.
    \item Обчислення енергії та електронної густини вимагає \textbf{інтегрування по зоні Бріллюена} --- через дискретну сітку $\mathbf{k}$-точок.
    \item Використовуються спеціальні \textbf{базисні набори}, адаптовані для періодичних систем (GTH).
\end{itemize}

%%% --------------------------------------------------------
%\subsection{Теорія Блоха та періодичні граничні умови}
%%% --------------------------------------------------------
%
%Основним математичним підґрунтям для опису періодичних систем є \textbf{теорія Блоха}. Відповідно до цієї теорії, хвильові функції електронів у періодичному потенціалі можуть бути записані як добуток двох функцій:
%
%\[
%\psi_{n\mathbf{k}}(\mathbf{r}) = e^{i \mathbf{k} \cdot \mathbf{r}} u_{n\mathbf{k}}(\mathbf{r}),
%\]
%де:
%\begin{itemize}
%    \item $\mathbf{k}$ --- вектор в зоні Бріллюена, що визначає хвильовий вектор електрону,
%    \item $u_{n\mathbf{k}}(\mathbf{r})$ --- періодична частина хвильової функції, яка має той самий період, що й ґратка кристала (тобто, $u_{n\mathbf{k}}(\mathbf{r}) = u_{n\mathbf{k}}(\mathbf{r} + \mathbf{R})$ для будь-якої транслокації $\mathbf{R}$ в межах ґратки).
%\end{itemize}
%
%Цей підхід дозволяє враховувати лише одну елементарну комірку та розглядати електронні стани через вектори хвильових чисел $\mathbf{k}$, що визначаються через дискретизацію зон Бріллюена.
%
%\section{Безкінечні періодичні кристали}
%
%Електронні властивості твердих тіл сильно залежать від їх складу та характеру хімічного зв'язку. Відповідно, способи опису різних кристалів суттєво відрізняються. Можна виділити два граничні випадки обчислення одноелектронних хвильових функцій у безкінечних періодичних кристалах.
%
%\subsection{1. Приближення сильної зв'язку}
%
%Інтерференція зовнішніх валентних атомних орбіталей відносно слабка, і атомна специфіка валентних електронів в значній мірі зберігається. У цьому випадку для опису хвильових функцій застосовується наближення ЛКАО, яке в фізиці твердого тіла називають наближенням сильної зв'язку. Це наближення добре описує молекулярні та атомні кристали.
%
%Ширина енергетичної смуги пропорційна інтегралу перекриття між сусідніми атомами, що відображає силу їх взаємодії. Зі збільшенням енергії валентного рівня просторове розширення його хвильової функції збільшується, а отже, збільшується ширина валентної зони. У металах валентна зона особливо широка, тому наближення сильної зв'язку має для них обмежену застосовність.
%
%
%\subsection{Приближення слабкої зв'язку}
%
%Якщо валентні електрони знаходяться на зовнішніх орбіталях, які слабо зв'язані з добре локалізованими атомними остовами, то інтерференція валентних хвильових функцій велика, і валентні електрони практично повністю втрачають свою атомну специфіку: вони колективізуються і набувають здатності переміщатися в кристалічній решітці. У такому випадку валентні електрони в першому наближенні можна вважати майже вільними і описувати плоскими хвилями, слабо модифікованими періодичним потенціалом решітки. Це наближення в фізиці твердого тіла називають наближенням слабкої зв'язку. Підхід в принципі застосовний лише до металів, однак на практиці він використовується в суттєво модифікованому вигляді, що значно розширює область його застосування (див. нижче).
%
%
%
%Приближення сильної зв'язку реалізується в сучасній практиці розрахунків за допомогою методів Хартрі—Фока і Кона—Шема в наближенні ЛКАО. Функції (6.5) використовуються як базисні, а періодичні в кристалічній решітці функції \( u_k(\mathbf{r}) \) найчастіше представляють собою локалізовані атомні орбіталі, які будуються з гаусових орбіталей. Як і в молекулярних розрахунках, застосовуються мінімальні та розширені базисні набори. Специфіка розрахунку кристала полягає в тому, що врахування періодичності еквівалентно включенню в базис атомних орбіталей у різних комірках, тому хороші результати досягаються вже при розрахунку з невеликою кількістю базисних функцій.
%
%Через врахування лише обміну та ігнорування миттєвої кулонівської кореляції електронів метод Хартрі—Фока переоцінює ширину забороненої зони кристала в 1,5—2 рази і непридатний для розрахунку Фермі-поверхні металів. У методі Кона—Шема ширина забороненої зони, навпаки, занижується.
%
%
%Приближення слабкої зв'язку реалізується за допомогою методу Кона—Шема, і різні його варіанти відрізняються способами апроксимації функцій \( u_k(\mathbf{r}) \) у (6.5). Розглянемо найбільш поширені підходи.
%
%Плоскі хвилі — точні власні функції задачі про однорідний електронний газ. Тому розкладання одноелектронних функцій \( f_k(\mathbf{r}) \) по плоских хвилях (ПХ) для металів, де розподіл валентної електронної густини майже однорідний, виглядає природно. Ці функції і були спочатку введені для розв'язання задач, пов'язаних з електронною структурою металів, а потім у модифікованому вигляді поширені на інші системи. ПХ ортогональні і не залежать від енергії атомних станів, а також від положення атомів у комірці. З математичної точки зору плоскі хвилі формують повний базисний набір, тому з збільшенням числа базисних функцій точність розв'язку одноелектронних рівнянь буде збільшуватися. Однак, щоб описати вузлову структуру хвильових функцій в області атомних остовів, потрібно більше \( 10^6 \) ПХ; через це збіжність методу дуже повільна.
%
%
%\subsection{Метод ортогоналізованих плоских хвиль (ОПВ)}
%
%В методі ортогоналізованих плоских хвиль (ОПВ або OPW) \cite{6.14} осциляції хвильових функцій остовів враховуються в базисі, по якому розкладаються одноелектронні функції. Для цього хвильові функції атомних остовів описуються блохівськими функціями, побудованими з локалізованих атомних орбіталей, до яких додаються плоскі хвилі, вимагаючи, щоб хвильові функції, які необхідно знайти, були ортогональні остовним хвильовим функціям. ОПВ також утворюють повну систему для розкладу функцій, що описують електрони провідності.
%
%
%\subsection{Метод приєднаних плоских хвиль (ППВ)}
%
%В іншому методі як базис для розкладу хвильових функцій використовуються так звані приєднані плоскі хвилі (ППВ або APW) \cite{6.15}. Ураховують, що вблизи атомних остовів потенціал близький до сферично симетричного, і тому потенціал кристала вибирають у вигляді центрованих на атомах неперекриваючихся сфер, а в просторі між сферами потенціал приймається рівним нулю (рис. 6.15, див. кольорову вкладку). Кристалічний потенціал, заданий таким чином, називається МТ-потенціалом*. Хвильові функції в методі ППВ визначаються за наступним рецептом. Всередині МТ-сфер, де потенціал сферично симетричний, будуються рішення атомного типу. В просторі між МТ-сферами хвильові функції представляють собою плоскі хвилі, які можна визначити для будь-якої енергії та будь-якого хвильового вектора. Коефіцієнти перед функціями в атомних областях підбираються таким чином, щоб на поверхні сфери хвильова функція гладко переходила в плоску хвилю (при цьому перша похідна ППВ на поверхні сфери має розрив). Кристалічні орбіталі представляються у вигляді лінійної комбінації ППВ, а коефіцієнти розкладу по ППВ для кожного значення хвильового вектора \( k \), що дозволено симетрією кристала, знаходять за допомогою варіаційного принципу. Для розрахунку електронної структури кристалів потрібно обмежене число ППВ.
%
%\subsection{Радіальні функції в межах атомних сфер та метод LAPW}
%
%Радіальні функції в межах атомних сфер \( u(l, r) \) в методі ППВ залежать від енергій \( E(l) \) (де \( l \) — кутове квантове число), які слід прийняти рівними енергіям відповідних зон \( e \). Оскільки саме ці енергії є метою розрахунку і заздалегідь невідомі, рішення ППВ виявляються нелінійними, що ускладнює процедуру самосогласованості при розв'язанні одноелектронних рівнянь. Щоб усунути цей недолік, було запропоновано \cite{6.16, 6.17} відмовитися від точного погодження енергій \( E(l) \) та \( e \) і застосувати розкладання функції \( u(l, r) \) в ряд Тейлора по \( e \) відносно \( E(l) \), обмежуючись лінійним членом:
%
%\[
%u(l, r) = u_0(l, r) + \frac{d u(l, r)}{d e} \Big|_{e=E(l)} (e - E(l)) + \cdots
%\]
%
%Це призвело до створення лінеаризованого методу приєднаних плоских хвиль (linearized augmented plane waves — LAPW), який вимагає безперервності на межі МТ-сфери не тільки для хвильових функцій, але й для їхніх похідних. Це особливо важливо, якщо для даного значення \( E(l) \) функція \( u(l, r) \) рівна нулю на межі МТ-сфери: у цьому випадку в методі ППВ відповідні локальні функції та плоскі хвилі не є погодженими, в той час як у методі LAPW цей недолік усувається. Часто достатньо використовувати лише один набір значень \( E(l) \), щоб описати валентну зону.
%
%\subsection{Метод FP-LAPW}
%
%Число базисних функцій на атом в методі LAPW складає близько 100. Рішення знаходять лінійним варіаційним методом, приводячи матрицю модельного гамільтоніана до діагонального вигляду; це є суттєвою перевагою LAPW порівняно з нелінійним методом ППВ. Крім того, базис LAPW забезпечує значну варіаційну свободу, що дозволяє легко ввести несферичне наближення для МТ-потенціалу. У цьому випадку метод називається повнопотенціальним лінеаризованим методом приєднаних плоских хвиль (full-potential linearized augmented plane-wave method — FP-LAPW) \cite{6.18—6.20}.
%
%\subsection{Метод LAPW+LO та LMTO}
%
%У деяких матеріалах є високоеенергетичні електронні остовні («півостовні») стани, які занадто делокалізовані, щоб їх можна було описати орбіталями атомного типу, що повністю знаходяться в межах МТ-сфери. Причому величини \( E(l) \), необхідні для їх опису, відповідають станам у валентній зоні. Щоб описати такі стани, в базис вводять додаткові локальні орбіталі (LO), повністю укладені всередині МТ-сфер (метод LAPW+LO): на межі МТ-сфери значення цих функцій та їх похідних дорівнюють нулю. Величини \( E(l) \) для цих функцій беруть рівними енергіям «півостовних» станів \cite{6.21, 6.22}. Базис у методі LAPW+LO розширюється порівняно з початковим методом LAPW, однак оскільки кількість плоских хвиль складає 50—100 на атом, пов'язане з цим збільшення часу розрахунку є незначним.
%
%Щоб досягти кращого опису в межах МТ-сфер, окрім базового LAPW-набору, вводять додатковий набір локальних орбіталей, незалежних від хвильових векторів. Значення локальних орбіталей на межах МТ-сфер дорівнюють нулю, а на похідні жодних обмежень не накладається. У цьому випадку кінетичну енергію електронів на поверхні сфери слід врахувати в гамільтоніані задачі (метод LAPW+LO) \cite{6.23}. Використання двох базисних наборів призводить до того, що збіжність розрахунку досягається на порядок швидше порівняно з LAPW \cite{6.24, 6.25}.
%
%МТ-приближення також використовується в іншій серії методів, в основі яких лежить лінійний метод МТ-орбіталей (linear muffin-tin orbital — LMTO). Потенціал у кристалі задається у формі МТ-сфер. Базисні функції будуються як лінійні комбінації блоховських функцій, компоненти яких у межах МТ-сфер описуються так само, як у методі ППВ, а між сферами апроксимуються функціями типу \( r^{-(l+1)} \); обидва рішення та їхні похідні на межі МТ-сфери зшиваються. Можна вибрати МТ-сфери, що проникають одна в одну таким чином, щоб їхній загальний об'єм став рівним об'єму елементарної комірки кристала; таке наближення називається LMTO-приближенням атомних сфер (atomic sphere approximation, LMTO-ASA).
%
%\subsection{Метод FP-LMTO та TB-LMTO}
%
%У повнопотенціальному методі LMTO (full potential linear muffin-tin orbital — FP-LMTO) потенціал та електронна густина в межах МТ-сфер є несферичними, що досягається їх розкладанням у ряд по кутових гармоніках з урахуванням локальної симетрії атомних позицій. Між МТ-сферами потенціал і густина розкладаються у ряди Фур'є, а потім міжатомні рішення зшиваються з атомними. З отриманих функцій будуються функції Блоха. Кількість базисних функцій на атом складає близько 15.
%
%LMTO-орбіталі можна використовувати як вихідні для побудови базису в методі сильної зв'язку (tight-binding LMTO, або TB-LMTO) \cite{6.26}. У цьому випадку базисні LMTO-орбіталі за допомогою унітарного перетворення трансформуються в добре локалізовані функції, з яких будуються кристалічні орбіталі.
%
%\subsection{Функції Ванньє та метод NMTO}
%
%Блохівські функції для будь-якої зони Бріллюєна завжди можна перетворити таким чином, що їх формальний вигляд буде таким же, як у випадку приближення сильної зв'язку. Функції, аналогічні атомним орбіталям, називаються функціями Ванньє. Для вузьких зон з сильною зв'язком функції Ванньє будуть добре локалізовані біля атомів. У модифікованому методі LMTO, методі NMTO (MTO N-го порядку) \cite{6.27-6.29}, функції, близькі по суті до функцій Ванньє, отримують шляхом трансформації мінімального базисного набору наступним чином. Потенціал вибирається у вигляді перекриваючихся МТ-сфер, рішення одноелектронних рівнянь в межах кожної сфери розкладаються в ряд по екранованим сферичним хвилям, апроксимуються поліномом, що описує залежність хвиль від енергії, а потім усувається залежність від енергії в (N + 1) точках за спеціальною схемою. При цьому частина функцій з базисного набору видаляється, а їх вплив враховується як ефект екранування. Суперпозиції результуючих функцій після симетричної ортогоналізації дають функції, подібні до функцій Ванньє, які добре локалізовані на атомах; їх кількість відповідає мінімальному базисному набору. Така досить складна процедура дозволяє в результаті аналізувати зонну картину в кристалі в рамках якісних хімічних уявлень і моделей, заснованих на взаємодії функцій, локалізованих на атомах \cite{6.30}.
%
%\subsection{Метод локалізованих орбіталей з повним потенціалом (FPLO)}
%
%Іншу схему використовує полнопотенціальний метод локалізованих орбіталей (full potential localized orbital — FPLO) \cite{6.31-6.32}. МТ-приближення для атомних потенціалів не застосовується; замість цього розмір потенціальної ями, в якій знаходиться атом, штучно обмежується. Кристалічні орбіталі в методі FPLO будуються з атомних орбіталей, кожна з яких апроксимована лише однією базисною функцією і визначається чисельно шляхом розв'язування одноелектронних рівнянь Кона—Шема. На атом припадає 30—40 базисних функцій. Недостатня гнучкість базису компенсується його оптимізацією в процесі розв'язку на кожному циклі самосогласування; в результаті вдається досягти результатів, точність яких не поступається методам LAPW.
%
%Для спрощення розв'язку одноелектронних рівнянь можна виключити з явного розгляду електрони атомних остовів, хвильові функції яких осцилюють, змінюючи знак у вузлах. Для цієї мети використовують валентне приближення: розподіл внутрішніх електронів в атомах, що входять до складу різних кристалів, вважається незмінним, а їх вплив зводиться до зміни ефективного заряду остову. Тоді замість кулонівського потенціалу, який описує електронно-ядрне взаємодія, можна використовувати більш слабкий потенціал в області остову — псевдопотенціал \cite{6.33-6.35}.
%
%Завдяки цьому хвильові функції валентних електронів замінюються псевдо-хвильовими функціями, які мало відрізняються від «правильних» функцій далеко від ядер і мають безвузлову природу в області атомних остовів (див. рис. 6.16). У результаті число ПВ у базисних функціях помітно зменшується (хоча й залишається на порядок більше, ніж у випадку локалізованих АО); окрім того, в псевдопотенціалі легко враховуються релятивістські ефекти.
%
%\subsection{Псевдопотенціал зберігаючий норму}
%
%Для забезпечення сшивання псевдопотенціального та атомного рішень варіаційної задачі псевдохвильова функція повинна бути неперервною і двічі диференційовною. Крім того, власні значення, що відповідають обом хвильовим функціям, повинні бути рівні, а заряди, зосереджені всередині сфери з радіусом сшивки, повинні збігатися. Потенціал, що задовольняє цим вимогам, називається псевдопотенціалом, що зберігає норму (norm-conserving pseudopotential) \cite{6.36}. Якщо кожній компоненті кутового моменту хвильової функції відповідає різний потенціал, то псевдопотенціал є нело-
%кальним.
%
%Псевдохвильова функція валентних електронів отримується після ортогоналізації рішень варіаційної задачі з псевдопотенціальним гамільтоніаном до хвильових функцій атомних остовів.
%
%\subsection{Псевдопотенціали та їх застосування}
%
%Псевдопотенціали можна знайти або з перших принципів \cite{6.37}, для чого необхідно знати хвильову функцію, отриману вирішенням одноелектронних рівнянь для заданої електронної конфігурації атома, або за допомогою підгонок параметрів, визначених з експериментальних даних (модельні та емпіричні псевдопотенціали). Завдання полягає в тому, щоб підібрати такий ефективний гамільтоніан, у якого електронний спектр збігається з точним в інтервалі енергій валентних електронів. Це можна зробити різними способами, тому псевдопотенціал визначається неоднозначно. Запропоновано кілька видів псевдопотенціалів \cite{6.39--6.45}, за допомогою яких успішно досліджуються енергетичні характеристики кристалів різного складу та природи.
%
%Схема, що в загальному вигляді охоплює основні наближення, які лежать в основі неемпіричних методів вирішення одноелектронних рівнянь (як методу Кона—Шэма, так і методу Хартрі—Фока) для структур з безкінечним періодичним потенціалом, показана на рис. 6.19. Ці методи обчислення кристалічних хвильових функцій використовують різні варіанти обмінно-кореляційних функціоналів, розглянуті раніше в главах 2 та 3; вони дозволяють також враховувати релятивістські ефекти в модельному гамільтоніані. У поєднанні з різноманітними можливостями вибору базисного набору, все це робить зазначені методи застосовними до широкого кола об'єктів — від молекулярних кристалів до інтерметалічних сполук, дозволяючи використовувати найбільш підходяще наближення в кожному конкретному випадку. Досяжна точність обчислення параметрів елементарних клітинок складає 0,5\%, а енергетичних характеристик — близько 0,1 еВ/атом.
%
%Ітераційна схема стандартного самосогласованого вирішення одноелектронних рівнянь з періодичним потенціалом наведена на рис. 6.20. Її відмінність від аналогічної схеми, що застосовується при обчисленнях молекул (див. рис. 3.4), полягає в тому, що рішення шукаються для всіх вибраних $k$-точок у першій зоні Бріллюена, а електронна щільність кристала отримується сумуванням одноелектронних функцій за дозволеними значеннями вектора $k$, а потім по смугах, зайнятих електронами (тобто до поверхні Фермі).
%
%Отже, ми з'ясували, що через періодичність кристалічної структури електрони в кристалі виявляються розподіленими по енергетичних зонах (смугах), які складаються з рівнів енергії, що відповідають хвильовим функціям з дозволеними симетрією хвильовими векторами. Властивості цих зон залежать від квантових чисел, що утворюють їх електронів, та від хімічної зв'язності між атомами.



%% --------------------------------------------------------
\subsection{Алгоритм розрахунків}
%% --------------------------------------------------------

Основні етапи розрахунку для періодичних систем:

\begin{enumerate}
    \item \textbf{Вибір типу потенціалу}: Спочатку вибирається тип псевдопотенціалу для заміни ефективними потенціалами остовних електронів.
    \item \textbf{Параметри ґратки}: Далі визначаються параметри елементарної комірки та її вектори, що задають періодичність.
    \item \textbf{Інтегрування по зоні Бріллюена}: Для точних результатів потрібна велика кількість $\mathbf{k}$-точок для дискретизації зони Бріллюена.
    \item \textbf{Спеціальні базисні набори}: Для систем із періодичними граничними умовами використовуються спеціалізовані базисні функції, що враховують симетрії періодичних структур, наприклад, \texttt{GTH} базисні функції, які адаптовані для кристалів і наноструктур.
\end{enumerate}

\subsection{Розв'язок рівнянь Шредінгера для періодичних систем}

Рішення рівнянь Шредінгера для періодичних систем, таких як кристали, поверхні та наноструктури, вимагає врахування періодичної природи потенціалу. Оскільки ґратка кристала нескінченна, розрахунки виконуються для лише однієї елементарної комірки, з подальшим застосуванням періодичних граничних умов, щоб отримати електронні стани для всієї структури.

Рівняння Шредінгера для таких систем можна записати як:

\[
H \psi_{n\mathbf{k}}(\mathbf{r}) = \epsilon_{n\mathbf{k}} \psi_{n\mathbf{k}}(\mathbf{r}),
\]
де $H$ --- гамільтоніан системи, а $\epsilon_{n\mathbf{k}}$ --- енергія електронного стану для хвильового вектора $\mathbf{k}$.


\section{Створення періодичної системи}

\subsubsection{Елементарна комірка}

Періодична система у \inlinecode{PySCF} створюється за допомогою класу \inlinecode{pyscf.pbc.gto.Cell}.
У цьому об’єкті описується склад комірки, геометрія ґратки, базис та псевдопотенціал.

\begin{minted}{python}
from pyscf.pbc import gto, scf, dft
import numpy as np

# Створення об'єкта комірки
cell = gto.Cell()

# Атоми всередині комірки (координати в Å або a.u.)
cell.atom = '''
C 0.0 0.0 0.0
C 0.8917 0.8917 0.8917
'''

# Вектори трансляції (решітка)
cell.a = '''
0.0    1.7834 1.7834
1.7834 0.0    1.7834
1.7834 1.7834 0.0
'''

# Базис і псевдопотенціал
cell.basis = 'gth-dzvp'
cell.pseudo = 'gth-pade'
cell.unit = 'Angstrom'

# Побудова об'єкта
cell.build()
\end{minted}

Після виконання \inlinecode{cell.build()} PySCF аналізує усі параметри, створює періодичну ґратку та перевіряє узгодженість розмірів.

\subsubsection{Типи ґраток}

Для зручності наведемо приклади задання ґраток різного типу.

\textbf{Кубічна решітка:}
\begin{minted}{python}
a = 5.0
cell.a = np.eye(3) * a
\end{minted}

\textbf{ГЦК (face-centered cubic):}
\begin{minted}{python}
a = 3.567  # наприклад, для Al
cell.a = [[0, a/2, a/2],
          [a/2, 0, a/2],
          [a/2, a/2, 0]]
\end{minted}

\textbf{Гексагональна (для графену):}
\begin{minted}{python}
a = 2.46
c = 20.0  # великий вакуумний зазор між шарами
cell.a = [[a, 0, 0],
          [a*np.cos(np.pi/3), a*np.sin(np.pi/3), 0],
          [0, 0, c]]
\end{minted}

\textbf{Порада:} для поверхонь і двовимірних матеріалів обов’язково робіть великий міжшаровий вакуум ($\ge 15$–20~Å), щоб запобігти взаємодії копій уздовж $z$.

\section{Базисні набори і псевдопотенціали}

У твердотільних розрахунках майже завжди застосовують \textbf{псевдопотенціали}.
Вони усувають остовні електрони з обчислень, залишаючи лише валентні, що суттєво зменшує кількість базисних функцій.

Найпоширеніші комбінації:
\begin{itemize}
    \item \inlinecode{basis='gth-dzvp', pseudo='gth-pade'} --- оптимальний баланс точності і швидкодії.
    \item \inlinecode{basis='gth-tzvp', pseudo='gth-pade'} --- вища точність, дорожче обчислення.
    \item \inlinecode{basis='gth-dzv', pseudo='gth-pade'} --- мінімальний базис для тестів.
\end{itemize}

GTH (Goedecker–Teter–Hutter) псевдопотенціали добре працюють як у розрахунках з плоскими хвилями, так і з ґаусовими базисами, які використовує PySCF.

\section{Сітка $k$-точок}

Оскільки властивості кристала періодичні, інтеграли беруться не по реальному простору, а по \textbf{зоні Бріллюена}.
Для цього її дискретизують на сітку $k$-точок --- аналог дискретизації простору в молекулярних розрахунках.

\begin{minted}{python}
# Генерація рівномірної сітки Монкхорста–Пака
kpts = cell.make_kpts([4, 4, 4])

# Для двовимірних систем
kpts = cell.make_kpts([8, 8, 1])

# Тільки Gamma-точка
kpts = cell.make_kpts([1, 1, 1])
\end{minted}

\textbf{Типові рекомендації:}
\begin{itemize}
    \item Метали: щільна сітка $8\times8\times8$ або густіша.
    \item Напівпровідники: $4\times4\times4$ часто достатньо.
    \item Великі суперкомірки (з дефектами): можна обмежитись лише $\Gamma$-точкою.
\end{itemize}

\section{SCF-розрахунки для твердих тіл}

\subsubsection{Метод Хартрі–Фока}

\begin{minted}{python}
from pyscf.pbc import scf

kmf = scf.KRHF(cell, kpts)
kmf.kernel()

print(f"Повна енергія: {kmf.e_tot:.6f} Ha")
print(f"Енергія на атом: {kmf.e_tot/len(cell.atom):.6f} Ha")
\end{minted}

K-точковий HF враховує періодичність хвильових функцій та дозволяє отримати енергію системи в самозгодженому полі.

\subsubsection{Метод DFT}

\begin{minted}{python}
from pyscf.pbc import dft

kmf = dft.KRKS(cell, kpts)
kmf.xc = 'PBE'
kmf.kernel()
\end{minted}

Функціонали для твердих тіл:
\begin{itemize}
    \item \texttt{LDA, VWN} --- простий, але часто недооцінює об’єм ґратки;
    \item \texttt{PBE} --- найпопулярніший GGA-функціонал для кристалів;
    \item \texttt{SCAN} --- мета-GGA, точніший, але повільніший;
    \item \texttt{B3LYP, HSE} --- гібридні, дорогі у розрахунках.
\end{itemize}

\section{Електронна структура і зони}

\begin{minted}{python}
mo_energy_kpts = kmf.mo_energy

for ik, kpt in enumerate(kpts):
    print(f"k-точка {ik}: {kpt}")
    print(f"Енергії орбіталей: {mo_energy_kpts[ik][:5]}")
\end{minted}

Таким чином отримуємо енергетичні рівні (зони) у кожній $k$-точці.
Для побудови зонної структури енергії можна з’єднати вздовж спеціальних шляхів у зоні Бріллюена.

\section{Типові приклади}

\subsubsection{Кристал кремнію}

\begin{minted}{python}
from pyscf.pbc import gto, dft
import numpy as np

cell = gto.Cell()
a = 5.431
cell.atom = '''
Si 0.00 0.00 0.00
Si 0.25 0.25 0.25
'''
cell.a = [[0, a/2, a/2],
          [a/2, 0, a/2],
          [a/2, a/2, 0]]
cell.basis = 'gth-dzvp'
cell.pseudo = 'gth-pade'
cell.unit = 'Angstrom'
cell.build()

kpts = cell.make_kpts([4, 4, 4])
kmf = dft.KRKS(cell, kpts)
kmf.xc = 'PBE'
e_tot = kmf.kernel()

print(f"Енергія на атом Si: {e_tot/2:.6f} Ha = {e_tot/2*27.2114:.3f} eV")
\end{minted}

\subsubsection{Графен (моношар)}

\begin{minted}{python}
cell = gto.Cell()
a = 2.46
c = 20.0
cell.atom = '''
C 0.0 0.0 0.0
C 0.0 1.42 0.0
'''
cell.a = [[a, 0, 0],
          [-a/2, a*np.sqrt(3)/2, 0],
          [0, 0, c]]
cell.basis = 'gth-dzvp'
cell.pseudo = 'gth-pade'
cell.dimension = 2
cell.build()

kpts = cell.make_kpts([8, 8, 1])
kmf = dft.KRKS(cell, kpts)
kmf.xc = 'PBE'
kmf.kernel()
\end{minted}

\subsubsection{Алюміній (метал)}

\begin{minted}{python}
cell = gto.Cell()
a = 4.05
cell.atom = 'Al 0 0 0'
cell.a = [[0, a/2, a/2],
          [a/2, 0, a/2],
          [a/2, a/2, 0]]
cell.basis = 'gth-dzvp'
cell.pseudo = 'gth-pade'
cell.build()

kpts = cell.make_kpts([8, 8, 8])
kmf = dft.KRKS(cell, kpts)
kmf.xc = 'PBE'

from pyscf.pbc.dft import numint
kmf.smearing_method = 'fermi'
kmf.sigma = 0.01

kmf.kernel()
\end{minted}

Для металів \textbf{обов’язково} використовуйте розмиття рівня Фермі (smearing), щоб уникнути числових осциляцій при частково заповнених зонах.

\section{Оптимізація геометрії}

\begin{minted}{python}
from pyscf.pbc.geomopt import optimize
mol_eq = optimize(kmf)
\end{minted}

PySCF може також взаємодіяти з пакетом \texttt{ASE} для гнучкої оптимізації:
\begin{minted}{python}
from ase.optimize import BFGS
from pyscf.pbc.tools import pyscf_ase

ase_atoms = pyscf_ase.to_ase_atoms(cell)
\end{minted}

\section{Аналіз результатів}

\subsubsection{Плотність станів (DOS)}

\begin{minted}{python}
import matplotlib.pyplot as plt

energies = np.concatenate(kmf.mo_energy)
plt.hist(energies, bins=100, density=True)
plt.xlabel('Енергія (Ha)')
plt.ylabel('Плотність станів')
plt.show()
\end{minted}

\subsubsection{Ширина забороненої зони}

\begin{minted}{python}
def get_band_gap(kmf):
    homo = max([e[e < kmf.mo_occ[ik][0]].max()
                for ik, e in enumerate(kmf.mo_energy)])
    lumo = min([e[e > kmf.mo_occ[ik][-1]].min()
                for ik, e in enumerate(kmf.mo_energy)])
    gap = (lumo - homo) * 27.2114
    return gap

gap = get_band_gap(kmf)
print(f"Ширина забороненої зони: {gap:.3f} eV")
\end{minted}

\section{Просунуті можливості}

\subsubsection{GW-поправка до зонної структури}

\begin{minted}{python}
from pyscf.pbc import gw

kmf = dft.KRKS(cell, kpts).run()
mygw = gw.KGWAC(kmf)
mygw.kernel()
\end{minted}

Метод \texttt{GW} дає квазичастинкові енергії, виправляючи недооцінений розрив зони у DFT.

\subsubsection{Поверхні та слеби}

\begin{minted}{python}
cell = gto.Cell()
cell.a = [[a, 0, 0],
          [0, a, 0],
          [0, 0, 30.0]]  # вакуум уздовж z
cell.dimension = 2
\end{minted}

\section{Практичні рекомендації}

\begin{enumerate}
    \item Починайте з малих систем і грубих сіток $k$-точок.
    \item Перевіряйте сходимість енергії при збільшенні сітки.
    \item Використовуйте \texttt{PBE} як базовий функціонал.
    \item Для молекул на поверхнях --- вакуум не менше 15~Å.
    \item Для металів --- обов’язково застосовуйте розмиття рівня Фермі.
\end{enumerate}

\section{Типові помилки}

\begin{itemize}
    \item \texttt{linear dependency in basis} --- базис занадто великий або лінійно залежний.
    → Використовуйте GTH-базиси.

    \item Повільна сходимість SCF --- погане початкове наближення.
    → Задайте \inlinecode{kmf.init\_guess = 'atom'}.

    \item Занадто велика енергія --- різні одиниці довжини.
    → Перевірте \inlinecode{cell.unit = 'Angstrom'}.

    \item Від’ємний band gap --- сітка $k$ занадто рідка.
    → Збільште кількість $k$-точок.
\end{itemize}

\section{Корисні ресурси}

\begin{itemize}
    \item Документація PySCF PBC: \url{https://pyscf.org/pbc.html}
    \item Приклади розрахунків: \url{https://github.com/pyscf/pyscf/tree/master/examples/pbc}
    \item База даних матеріалів: \url{https://materialsproject.org}
\end{itemize}
